%; whizzy chapter -dvi
% -initex iniptex -latex platex -format platex -bibtex jbibtex -fmt fmt
% 以上 whizzytex を使用する場合の設定。
 
%     Tokyo Debian Meeting resources
%     Copyright (C) 2012 Junichi Uekawa
%     Copyright (C) 2011 Nobuhiro Iwamatsu

%     This program is free software; you can redistribute it and/or modify
%     it under the terms of the GNU General Public License as published by
%     the Free Software Foundation; either version 2 of the License, or
%     (at your option) any later version.

%     This program is distributed in the hope that it will be useful,
%     but WITHOUT ANY WARRANTY; without even the implied warranty of
%     MERCHANTABILITY or FITNESS FOR A PARTICULAR PURPOSE.  See the
%     GNU General Public License for more details.

%     You should have received a copy of the GNU General Public License
%     along with this program; if not, write to the Free Software
%     Foundation, Inc., 51 Franklin St, Fifth Floor, Boston, MA  02110-1301 USA

%  preview (shell-command (concat "evince " (replace-regexp-in-string "tex$" "pdf"(buffer-file-name)) "&"))

%%ここからヘッダ開始。

\documentclass[mingoth,a4paper]{jsarticle}
\usepackage{monthlyreport}
% 日付を定義する、毎月変わります。
\newcommand{\debmtgyear}{2014}
\newcommand{\debmtgmonth}{06}
\newcommand{\debmtgdate}{14}
% started from zero:
% (let ((year 2013) (month 7)) (+ (* (- year 2005) 12) month -1))
\newcommand{\debmtgnumber}{114}

% tikz picture の為のマクロ設定
\usepackage[dvipdfmx]{graphicx}
\usepackage{tikz}

\begin{document}

\begin{titlepage}
\thispagestyle{empty}
% タイトルページ:編集必要な部分は最初のマクロに飛ばすこと

\vspace*{-2cm}
第\debmtgnumber{}回 東京エリア Debian 勉強会資料\\
\hspace*{-2cm}
\includegraphics{image2012-natsu/dotdeb.pdf}\\
\hfill{}\debmtgyear{}年\debmtgmonth{}月\debmtgdate{}日

% ここはアップデートすること
% 全角文字にしないとフォントのサイズが合わないので注意
\rotatebox{10}{\fontsize{32}{32} {\gt 特集：GPG 秘密鍵取り扱い方法の提案}}\\

\vspace*{-2cm}
\hfill{}\includegraphics[height=6cm]{image200502/openlogo-nd.eps}
\end{titlepage}

\newpage

\begin{minipage}[b]{0.2\hsize}
 \definecolor{titleback}{gray}{0.9}
 \colorbox{titleback}{\rotatebox{90}{\fontsize{80}{80} {\gt デビアン勉強会} }}
\end{minipage}
\begin{minipage}[b]{0.8\hsize}
\hrule
\vspace{2mm}
\hrule
\begin{multicols}{2}
\tableofcontents
\end{multicols}
\vspace{2mm}
\hrule
\end{minipage}

\dancersection{事前課題}{野島 貴英}

今回の事前課題は以下です:
\begin{enumerate}
 \item 本日、何の作業をやるかを宣言ください。
\end{enumerate}
この課題に対して提出いただいた内容は以下です。
\begin{multicols}{2}
{\small
 \begin{prework}{ wbcchsyn }
発表予定のツールをブラッシュアップする。\\
\url{https://github.com/wbcchsyn/dVault\_prototype}
フィードバックがあれば、取り込み。\\
無ければ、ドキュメント作成とか自分で気がついている部分の改修とか。
\end{prework}

\begin{prework}{ 野島　貴英 }
\begin{itemize}
\item Debian ARM64まわりを調べる。
\item OSS化したsoftetherをdebianパッケージとしてどうするか考える
\item Debian Gnu/Hurdネタを調べる
\end{itemize}

 まあ、将来の東京エリアDebian勉強会ネタの仕込み作業となりますなぁ。
\end{prework}

\begin{prework}{ yyuu }

 先々月に続いて、pyenv というスクリプトの deb 化のための作業を行なおうと考えています。

 \url{https://github.com/yyuu/pyenv}

 可能であれば pyenv-virtualenv プラグインも deb 化したいです。

 \url{https://github.com/yyuu/pyenv-virtualenv}

\end{prework}

\begin{prework}{ zinrai }

direnvのdebパッケージ化

\url{https://github.com/zimbatm/direnv}

\end{prework}

\begin{prework}{ koedoyoshida }

あんどきゅめんてっどでびあんの編集作業

\end{prework}

\begin{prework}{ regonn(Kenta Tanoue) }

AWS debian環境で、
Rails+publifyを使って自分のブログを作ります。（作業工程を記事にしたい）
インストールの隙間時間でDDTSSの作業もできたらいいな。
\end{prework}

\begin{prework}{ なかおけいすけ }

7月の勉強会で発表せよとご指名いただいたので、その準備をします。

\end{prework}

\begin{prework}{ 野首(@knok) }

\begin{itemize}
\item KAKASIのバグフィックス
\item gnu.org翻訳
\item ECS Livaでのdebian kernelでのmmc関連修正
\end{itemize}

\end{prework}

}
\end{multicols}

\dancersection{Debian Trivia Quiz}{野島 貴英}

ところで、みなさん Debian 関連の話題においついていますか？Debian関連の話
題はメーリングリストをよんでいると追跡できます。ただよんでいるだけではは
りあいがないので、理解度のテストをします。特に一人だけでは意味がわからな
いところもあるかも知れません。みんなで一緒に読んでみましょう。

今回の出題範囲は\url{debian-devel-announce@lists.debian.org} や \url{debian-devel@lists.debian.org}に投稿された
内容などからです。

\begin{multicols}{2}
%; whizzy-master ../debianmeetingresume201311.tex
% 以上の設定をしているため、このファイルで M-x whizzytex すると、whizzytexが利用できます。
%

\santaku
{MATE　desktop環境が先日Debianのパッケージに入りました。対象のMATEのバージョンは？}
{1.7}
{1.8}
{1.9}
{B}
{Gnome 2.28ベース以降のGnomeからフォークしたデスクトップ環境とのことです。読み方は「マテ」（マテ茶のマテ）と読むとのこと。}

\santaku
{ARM64アーキテクチャへの対応の協力者募集が行われました。残念ながらDebianは動きませんが、民生品で手に入るARM64搭載の機材は次のどれでしょう？}
{iphone 5s}
{NEXUS 7}
{OpenBlocks A7}
{A}
{iphone 5sから搭載されているA7 SoCはARMv8搭載でARM64動作も出来るそうです。ARM64を搭載した民生品はこれが最初かも。
一方、DebianのARM64ポーティングでは全く別の開発用向けに提供されている機材が利用されています。
我々がDebianの動くARM64機材を民生品で手に出来る日は何時ー（涙）？}

\santaku
{edos.debian.netがqa.debian.org/doseで復活しました。ところで、このサイト何するサイト？}
{各パッケージの依存関係がお互いに満たされている状態かをチェック}
{各パッケージがどれだけupstreamから乖離しているかチェック}
{各パッケージがどれだけ人気があるかをチェック}
{A}
{dose-*パッケージに含まれるチェックコマンドを使い、各アーキテクチャのdebianリポジトリを毎日チェックして、依存関係が満たされていないパッケージをチェックした結果を教えるサイトとなります。膨大な量のパッケージの依存関係を迅速にチェック出来るところが凄い点です。仕組みについてIEEEの論文まで出ているようです。参考：http://hal.archives-ouvertes.fr/ docs/00/14/95/66/PDF/ase.pdf}

\santaku
{先日Debian GNU/Hurdの中間報告がなされました。initが変更になったそうですが、どれになったでしょうか？}
{Hurd独自実装のinit}
{systemd}
{Sysvinit}
{C}
{去年のGSoCの成果がそのまま開発継続し、遂にSysvinitにinitが変更されたそうです。このおかげで、Network、halt/shutdown時の動作、ブート時のfilesystemのマウント方式がやっとDebian流になったようです。}

\santaku
{DEP-12がPTSによりサポートされたそうです。ところで、DEP-12の内容って何？}
{CIツールによる自動テストの為の定義ファイルの提案}
{upstreamに関する様々な情報を記載したメタデータの提案}
{copyrightファイルを機械処理がしやすいフォーマットにする提案}
{B}
{DEP-12はソースパッケージ中の''debian/upstream/metadata''等のファイルをに用意し、upstreamに関する様々な情報を記載しようという提案です。
 一方、「CIツールによる自動テストの為の...」はDEP-8であり、こちらは''debian/tests/''以下のファイルを使い、ci.debian.netでパッケージ開発について継続的インテグレーションを実施するものとなります。先日紹介のアナウンスありましたね。}

\santaku
{dputする前にはduckで必ずテストしてほしいとのこと。ところでduckって何？}
{debian/以下のファイルに含まれるURL/メールアドレスの存在をチェックするツール}
{状況・様子から名前を断定するツール}
{依存関係をチェックするツール}
{A}
{debian/controlファイル、debian/upstream関係のファイルに記載されているURL,メールアドレスのドメインパート(例:hoge@fuga.comのfuga.com部分)が実際に存在するか/アクセス可能かを調べてくれるツールがduckとなります。}

\santaku
{mentors.debian.netでレビュー待ちのパッケージはいくつある？}
{40〜50}
{50〜60}
{70以上}
{C}
{東京エリアDebian勉強会へいらっしゃるような皆様は、ぜひdebian-devel@debian.or.jpなどでDDにmentor役をお願いする、あるいはdebianの開発チームでパッケージがメンテナンスをされている場合はチームの誰かにmentor役をお願いするのがパッケージアップロードの道として早いです。不定期に開かれるDebian Hack Cafeで面前で直接頼むという手もあります。}

\end{multicols}

\dancersection{最近のDebian関連のミーティング報告}{野島 貴英}

\subsection{東京エリアDebian勉強会113回目報告}

 東京エリアDebian勉強会113回目は(株)スクウェア・エニックスさんで開催されました。

 8名の参加者がありました。また、野島さんより、仮想コンテナのフレームワークである
docker.ioについての発表が、実演ベースで行われました。
 
 また、残った時間は各自でもくもく作業をして、結果を発表しました。

　宴会は、会場近くの「南国亭新宿店」にて行いました。

% % (query-replace-regexp "<.*?>" "")
% % (query-replace-regexp "^[	 ]\+" "")

%-------------------------------------------------------------------------------
\dancersection{GPG 秘密鍵取り扱い方法の提案}{吉田 晋}
%-------------------------------------------------------------------------------
\index{gpg-private-keys-store-and-maniplate}

\subsection{はじめに}

皆さん、GPG のキーサインパーティーに参加した事ありますか？
私は、昨年の大統一 Debian 勉強会 2013 で初めて参加し、GPG の秘密鍵も、その時初めて作りました。
ところで、この秘密鍵はどのように保管すれば良いのでしょうか？

鍵はいつでも再発行出来るとはいえ色々と手間がかかりますし、古い鍵で暗号化したファイルを複合化出来なくなると困ります。
記録メディアの破損による鍵のロストには要注意です。
一般的に、記録メディアの破損に備えるには複数のメディアを複数の場所で保存する事は良い方法です。

一方、最低限のマナーとして盗難にも気をつける必要が有るでしょう。
厳重に管理するには、バックアップの数は少ないほど有利です。

さて、メディア破損に備えて冗長性を持たせる方法と盗難に備えて厳重に管理する方法は真っ向から矛盾してしまいました。
結局、どうすれば良いのでしょう？

\subsection{アイディアのベースは RAID}

例えば 3 本のディスクを用いた RAID5 のディスクアレイでは、2 個のディスクにデータを分散して記録しながら残りの 1 個のディスクにパリティデータを保存します。
同じ要領で秘密鍵を 2 個のファイル + 1 個のパリティファイル = 計 3 個のファイルを PC、USB メモリ（肌身離さず）、DVD（金庫）の 3 カ所で保存すれば、
\begin{itemize}
\item 普段は、USB メモリと PC を使って、都度鍵ファイルを復元して使用

  （使用後は鍵ファイル削除。削除忘れ防止のため、tmpfs 上で作ると良いかも。）
\item PC や USB メモリが盗まれても、しばらくは安全
\item ディスクのクラッシュや USB メモリの破損時には金庫の DVD を使って復元
\end{itemize}
という冗長性と秘匿性を兼ね備えた状態に成ります。
完全な鍵を世界中のどこにも保存する必要が無いところがポイント高！

\begin{figure}[H]
\begin{center}
  \begin{tikzpicture} [x=1em, y=1ex]
    \draw (-7, 0) rectangle (7, -24.5);
    \draw (0, -3) node {秘密鍵};
    \draw [red] (-3.8, -7.5) rectangle (3.8, -13.5);
    \draw (0, -9) node {1 0 0 0 1 1 ...};
    \draw (0, -12) node {......};
    \draw [blue] (-3.8, -17) rectangle (3.8, -23);
    \draw (0, -18.5) node {0 0 1 1 1 0 ...};
    \draw (0, -21.5) node {......};

    \draw [red] (-23, -30.5) rectangle (-9, -42.5);
    \draw (-16, -33) node {PC のディスクで保持)};
    \draw (-16, -36) node {1 0 0 0 1 1 ...};
    \draw (-16, -39) node {......};

    \draw [blue] (-7, -30.5) rectangle (7, -42.5);
    \draw (0, -33) node {USB メモリで肌身話さず};
    \draw (0, -36) node {0 0 1 1 1 0 ...};
    \draw (0, -39) node {......};

    \draw [green] (9, -30.5) rectangle (23, -42.5);
    \draw (16, -33) node {金庫で保管(parity)};
    \draw (16, -36) node {1 0 1 1 0 1 ...};
    \draw (16, -39) node {......};
  \end{tikzpicture}
\caption{分散管理イメージ図}
\end{center}
\end{figure}

ちなみに、分散するファイルの数を 2 倍にしたらどうなるでしょう？
秘密鍵を 4 個のファイル + 2 個 = 計 6 個のパリティファイルに分散すれば、パリティ増による冗長性の向上と、ファイル 1 個あたりの情報量減による秘匿性の向上を両立可能です。

\subsection{プロトタイプ作ってみた}

プロトタイプは Python 2.7 で作成し、 Debian Sid で動作確認を行いました。

dVault (distribute vault) と命名
\vspace{1cm}

\url{https://github.com/wbcchsyn/dVault_prototype}

（src/dvault が本体）
\vspace{1cm}

\subsubsection{注意}

あくまでプロトタイプなので、絶対に自分の秘密鍵に対して使わないで下さい。

例えば、
\begin{itemize}
\item 作成したファイルのパーミッションとか、全然気にしてません。
\item 問答無用でファイルの上書きをします。

  秘密鍵を上書きしながら途中で失敗すると、秘密鍵が失われます。
\end{itemize}
\vspace{1cm}

\subsubsection{使用方法}

origin というファイルを split0, split1, parity の 3 個のファイルに分割する場合

\begin{commandline}
$ dvault split -u origin -s split0 -s split1 -s parity
\end{commandline}

split0, parity の 2 個のファイルから origin を復活させる場合

\begin{commandline}
$ dvault unite -u origin -s split0 -s split1 -s parity
\end{commandline}
\vspace{1cm}

\subsubsection{仕様}
\begin{itemize}
\item パリティは 1 個で固定
\item 分散した各ファイルの最初にはメタ情報を json で保存。

  その後に 1 行空行を入れて、実際のデータを記入

  とりあえず、メタ情報は以下の 3 個 \begin{itemize}
  \item united\_length 元ファイルの長さ
  \item split\_count 何個に分割したか。（最低 2。）

    この数とパリティファイル 1 個、最低 3 個のファイルが出来る
  \item index 何番目のファイルか。（0 始まり。）

    index = split\_count の時は、パリティファイル
  \end{itemize}
\item 元ファイルを 1 byte ずつ分割。

  united\_length が split\_count で割り切れない場合は 0 でフィリング
\item 最後のファイル書き出し時以外は全てオンメモリで処理。

  一時ファイルは作りません
\end{itemize}
split\_count = 2 の場合、奇数目 byte を index = 0 のファイルに、偶数目 byte を index = 1 のファイルに、パリティ情報を index = 2 のファイルに保存

\subsection{実演}
\subsubsection{分割}

中身が \verb|abcedf\n| の 7 byte のファイルを 2 個 + parity の計 3 個に分割してみる
\vspace{1cm}

\begin{commandline}
$ hexdump -C origin

00000000  61 62 63 64 65 66 0a                              |abcdef.|
00000007

$ dvault split -u origin -s split0 -s split1 -s parity

$ hexdump -C split0

00000000  7b 22 69 6e 64 65 78 22  3a 20 30 2c 20 22 75 6e  |{"index": 0, "un|
00000010  69 74 65 64 5f 6c 65 6e  67 74 68 22 3a 20 37 2c  |ited_length": 7,|
00000020  20 22 73 70 6c 69 74 5f  63 6f 75 6e 74 22 3a 20  | "split_count": |
00000030  32 7d 0a 0a 61 63 65 0a                           |2}..ace.|
00000038

$ hexdump -C split1

00000000  7b 22 69 6e 64 65 78 22  3a 20 31 2c 20 22 75 6e  |{"index": 1, "un|
00000010  69 74 65 64 5f 6c 65 6e  67 74 68 22 3a 20 37 2c  |ited_length": 7,|
00000020  20 22 73 70 6c 69 74 5f  63 6f 75 6e 74 22 3a 20  | "split_count": |
00000030  32 7d 0a 0a 62 64 66 00                           |2}..bdf.|
00000038

$ hexdump -C parity

00000000  7b 22 69 6e 64 65 78 22  3a 20 32 2c 20 22 75 6e  |{"index": 2, "un|
00000010  69 74 65 64 5f 6c 65 6e  67 74 68 22 3a 20 37 2c  |ited_length": 7,|
00000020  20 22 73 70 6c 69 74 5f  63 6f 75 6e 74 22 3a 20  | "split_count": |
00000030  32 7d 0a 0a 03 07 03 0a                           |2}......|
00000038
\end{commandline}

\subsubsection{リストア}

フィアル split1、origin を無くしたとして、split0、parity から origin を復活させる
\vspace{1cm}

\begin{commandline}
$ rm origin split1

$ dvault unite -u origin -s split0 -s parity

$ hexdump -C origin

00000000  61 62 63 64 65 66 0a                              |abcdef.|
00000007
\end{commandline}

\subsection{悩み}
\begin{itemize}
\item Python で実装したが、果たして良かったのか？

  Debian の最小構成で Python が入っていない

  パッケージ化を目指すなら、C で再実装した方が良いかも
\item メタ情報を json で保存して良かったのだろうか？

  今後メタ情報が増えた場合に、文字列と数字の区別が出来たら便利だと思って json にした。

  しかし、C で実装する場合は json のパースは面倒。

  何より、現状では json の中で改行 x 2 が入ったらバグるので何とかする必要あり。
\item オンメモりの処理にどこまでこだわるか？

  OS 上のファイルとしては存在しなくとも、ディスクに少しでもデータが残ると盗難時に解析される恐れがあるのでオンメモリの処理にこだわった。

  でも、よく考えると SWAP したら意味が無いし、最初に秘密鍵を生成する際にディスク上に一時保存したら終わり。

  大容量のファイルを分割する事に備えて、名無しの一時ファイル（開いた直後に OS で削除）くらいは使っても良いかも。
\item 複数パリティに対応した方が良いか？
\end{itemize}

\subsection{今後の予定}

自分以外の人が使ってくれるなら、メンテナンスを行うつもりです。

（ドキュメントも作成する必要あり）

「使ってみたい」という人がいらしたら、教えて下さい。

%-------------------------------------------------------------------------------
\dancersection{会場での無線LANのつなぎ方}{野島 貴英}
%-------------------------------------------------------------------------------
 \subsection{はじめに}

　今回試験として、会場側でフィルタ無しのグローバル回線を用意しました。
ただ、会場側のセキュリティポリシーにより、
wpa-psk AES hidden SSIDという方式での提供となります。

　以下にDebianマシンでの接続方法を記載します。

 また、自分の環境では違うやり方でつながったという方は、野島まで
教えて下さい。こちらでもノウハウとして溜めていく予定です。

 \subsection{wpasupplicant及び/etc/network/interfacesを利用の場合}

 もっとも良いマニュアルは、/usr/share/doc/wpasupplicant/README.Debian.gz
となります。困った場合はこちらも合わせてご参照下さい。

　以下に/etc/network/interfacesの定義について会場の例を記載します。

\begin{commandline}  
$ sudo aptitude install wpasupplicant
# hidden ssidの元では必ず ap-scan 1,scan-ssid 1を指定する事。
# 参考：http://bugs.debian.org/358137
$ sudo vi /etc/network/interfaces
-----以下のエントリを追記ここから----------
iface wlan_tokyodebian inet dhcp
     wpa-ssid <<会場のSSID>>
     wpa-psk  <<会場のパスワード>>
     wpa-ap-scan 1
     wpa-scan-ssid 1
     
-----以下のエントリを追記ここまで----------
#　無線LANを有効にする。
$ sudo ifup wlan0=wlan_tokyodebian
#　無線LANを無効にする。
$ sudo ifdown wlan0
\end{commandline}
%$
 また、ハマってしまった時のデバッグ方法は、/usr/share/doc/wpasupplicant/README.Debian.gz中の''4. Trubleshooting''の章が便利です。

 \subsection{その他の無線LAN用パッケージを利用の場合}

　すみません、自分が情報を持たないため、現場で教えて下さい。

\printindex

\cleartooddpage

\vspace*{15cm}
\hrule
\vspace{2mm}
\includegraphics[width=2cm]{image200502/openlogo-nd.eps}
\noindent \Large \bf Debian 勉強会資料\\
\noindent \normalfont \debmtgyear{}年\debmtgmonth{}月\debmtgdate{}日 \hspace{5mm}  初版第1刷発行\\
\noindent \normalfont 東京エリア Debian 勉強会 （編集・印刷・発行）\\
\hrule

\end{document}
